\documentclass[11pt, a4paper]{ctexart}

\usepackage{experiment_style}

\begin{document}

% =========================================
% P61. 2
% =========================================
\begin{problem}{P61. 2}
    试用 MATLAB 软件编程实现追赶法求解三对角方程组的算法，并考虑如下梯形电阻电路问题。

    设电路中的各个电流 $\{i_1, i_2, \dots, i_8\}$ 满足线性方程组
    \[
        \begin{aligned}
            2 i_1 - 2 i_2             &= V / R, \\
            -2 i_1 + 5 i_2 - 2 i_3    &= 0, \\
            -2 i_2 + 5 i_3 - 2 i_4    &= 0, \\
            -2 i_3 + 5 i_4 - 2 i_5    &= 0, \\
            -2 i_4 + 5 i_5 - 2 i_6    &= 0, \\
            -2 i_5 + 5 i_6 - 2 i_7    &= 0, \\
            -2 i_6 + 5 i_7 - 2 i_8    &= 0, \\
            -2 i_7 + 5 i_8            &= 0.
        \end{aligned}
    \]
    设 $V = 220\,\text{V}, R = 27\,\Omega$，求各段电路的电流量。
\end{problem}

\begin{solution}
    \textbf{实验内容、步骤及结果}

    先编写追赶法函数 \texttt{chap-2/thomas\_algorithm.m}：

\begin{lstlisting}[language=Matlab]
function x = thomas_algorithm(a, b, c, d)
    % a: sub-diagonal (length n-1)
    % b: main diagonal (length n)
    % c: super-diagonal (length n-1)
    % d: right-hand side (length n)
    
    n = length(b);
    
    % Forward elimination
    for i = 2:n
        w = a(i-1) / b(i-1);
        b(i) = b(i) - w * c(i-1);
        d(i) = d(i) - w * d(i-1);
    end
    
    % Back substitution
    x = zeros(n, 1);
    x(n) = d(n) / b(n);
    for i = n-1:-1:1
        x(i) = (d(i) - c(i) * x(i+1)) / b(i);
    end
end
\end{lstlisting}

    然后建立本题对应的三对角方程组并调用追赶法程序 \texttt{chap-2/Q2.m}：

\begin{lstlisting}[language=Matlab]
% Define input parameters
a = -2;
a = repmat(a, 1, 7);
b = [2, 5, 5, 5, 5, 5, 5, 5];
c = a;
V = 220;
R = 27;
d = [V/R, 0, 0, 0, 0, 0, 0, 0];

% Call the function
x = thomas_algorithm(a, b, c, d);

% Display the output
disp(x);
\end{lstlisting}

    \textbf{实验结果分析}

\begin{lstlisting}
>> Q2
    8.1478
    4.0737
    2.0365
    1.0175
    0.5073
    0.2506
    0.1194
    0.0477
\end{lstlisting}

    因此，各段电流近似为
    \[
        (i_1, i_2, \dots, i_8)^{\mathrm{T}}
        \approx
        (8.1478, 4.0737, 2.0365, 1.0175, 0.5073, 0.2506, 0.1194, 0.0477)^{\mathrm{T}}.
    \]
    可见电流沿梯形电阻电路逐段递减。
\end{solution}


% =========================================
% P61. 3
% =========================================
\begin{problem}{P61. 3}
    方程组的性态和矩阵条件数实验。设有线性方程组 $Ax = b$，其中系数矩阵
    $A = (a_{ij})_{n \times n}$ 的元素分别为
    \[
        a_{ij} = x_i^{j-1}
        \quad
        \left(x_i = 1 + 0.1 i,\ i, j = 1, 2, \dots, n \right)
    \]
    或
    \[
        a_{ij} = \frac{1}{i + j - 1}
        \quad
        (i, j = 1, 2, \dots, n),
    \]
    右端向量
    \[
        \boldsymbol{b} =
        \left(
            \sum_{j=1}^n a_{1j},
            \sum_{j=1}^n a_{2j},
            \dots,
            \sum_{j=1}^n a_{nj}
        \right)^{\mathrm{T}}.
    \]
    利用 MATLAB 中的库函数完成下列实验：
    \begin{enumerate}[1.]
        \item 取 $n = 5, 10, 20$，分别求出系数矩阵的 2-条件数，判别它们是否为病态阵，并说明随着 $n$ 的增大矩阵性态如何变化。
        \item 分别取 $n = 5, 10, 20$，解上述两个线性方程组 $Ax = b$。将求得的解与精确解 $x = (1, 1, \dots, 1)^{\mathrm{T}}$ 作比较，并说明实验现象。
        \item 取 $n = 10$，对系数矩阵中的 $a_{22}$ 和 $a_{nn}$ 增加一个扰动 $10^{-8}$，求解方程组 $(A + \delta A)x = b$，观察解的变化。
    \end{enumerate}
\end{problem}

\begin{solution}
    \textbf{实验内容、步骤及结果}

    为便于统一生成两类线性方程组，先编写辅助函数 \texttt{chap-2/generateLinearSystem.m}：

\begin{lstlisting}[language=Matlab]
function [A, b] = generateLinearSystem(n, matrixType)
    % This function generates the coefficient matrix A and the vector b for
    % the linear system Ax = b. The parameter n specifies the size of the
    % matrix, and matrixType specifies the type of matrix to generate.
    %
    % matrixType can be:
    % 'polynomial' - A is generated using a_{i j} = x_i^{j-1} where x_i = 1 + 0.1 * i
    % 'hilbert'    - A is generated using a_{i j} = 1 / (i + j - 1)

    if strcmp(matrixType, 'polynomial')
        % Generate polynomial matrix
        A = zeros(n);
        for i = 1:n
            x_i = 1 + 0.1 * i;
            for j = 1:n
                A(i, j) = x_i^(j-1);
            end
        end
    elseif strcmp(matrixType, 'hilbert')
        % Generate Hilbert matrix
        A = zeros(n);
        for i = 1:n
            for j = 1:n
                A(i, j) = 1 / (i + j - 1);
            end
        end
    else
        error('Invalid matrix type. Use ''polynomial'' or ''hilbert''.');
    end

    % Compute vector b
    b = sum(A, 2);
end
\end{lstlisting}

    \textbf{(1) 计算条件数}

    条件数实验程序 \texttt{chap-2/Q31.m}：

\begin{lstlisting}[language=Matlab]
% Define the values of n
n_values = [5, 10, 20];
matrixType = 'polynomial'; % or 'hilbert'

% Loop through each value of n
for n = n_values
    % Generate the linear system
    [A, b] = generateLinearSystem(n, matrixType);
    
    % Compute the 2-norm condition number
    cond_num = cond(A, 2);
    
    % Determine if the matrix is ill-conditioned
    if cond_num > 1e10 % Cond(A) >> 1
        result = 'Ill-conditioned';
    else
        result = 'Well-conditioned';
    end
    
    % Display the condition number and the result
    fprintf('%s, n = %d, 2-norm condition number = %.2e, Result: %s\n', matrixType, n, cond_num, result);
end
\end{lstlisting}

    运行结果如下：

\begin{lstlisting}
>> Q31
polynomial, n = 5, 2-norm condition number = 5.36e+05, Result: Well-conditioned
polynomial, n = 10, 2-norm condition number = 8.68e+11, Result: Ill-conditioned
polynomial, n = 20, 2-norm condition number = 6.07e+21, Result: Ill-conditioned
>> Q31
hilbert, n = 5, 2-norm condition number = 4.77e+05, Result: Well-conditioned
hilbert, n = 10, 2-norm condition number = 1.60e+13, Result: Ill-conditioned
hilbert, n = 20, 2-norm condition number = 2.11e+18, Result: Ill-conditioned
\end{lstlisting}

    由结果可见，随着 $n$ 增大，两类矩阵的 2-条件数迅速增大，矩阵病态程度显著增强。

    \textbf{(2) 比较数值解与精确解}

    求解线性方程组并计算相对误差的程序 \texttt{chap-2/Q32.m}：

\begin{lstlisting}[language=Matlab]
% Define the values of n
n_values = [5, 10, 20];
matrixTypes = {'polynomial', 'hilbert'};

% Exact solution
x_exact = @(n) ones(n, 1);

% Loop through each matrix type and value of n
for matrixType = matrixTypes
    for n = n_values
        % Generate the linear system
        [A, b] = generateLinearSystem(n, matrixType{1});
        
        % Solve the linear system
        x_approx = A \ b;
        
        % Compute the relative error using the infinity norm
        rel_error = norm(x_approx - x_exact(n), inf) / norm(x_exact(n), inf);
        
        % Display the relative error
        fprintf('%s, n = %d, Relative Error (inf-norm) = %.2e\n', matrixType{1}, n, rel_error);
    end
end
\end{lstlisting}

    运行结果如下：

\begin{lstlisting}
>> Q32
polynomial, n = 5, Relative Error (inf-norm) = 2.89e-11
polynomial, n = 10, Relative Error (inf-norm) = 7.21e-05
polynomial, n = 20, Relative Error (inf-norm) = 3.50e+05
hilbert, n = 5, Relative Error (inf-norm) = 6.17e-13
hilbert, n = 10, Relative Error (inf-norm) = 3.18e-04
hilbert, n = 20, Relative Error (inf-norm) = 2.64e+01

Warning: Matrix is close to singular or badly scaled. Results may be inaccurate. RCOND =  3.713214e-23.
> In Q32 (line 15): x_approx = A \ b;
Warning: Matrix is close to singular or badly scaled. Results may be inaccurate. RCOND =  1.276108e-19.
> In Q32 (line 15): x_approx = A \ b;
\end{lstlisting}

    随着 $n$ 增大，系数矩阵的病态程度进一步提升，造成解的相对误差明显增大；当矩阵高度病态时，即使理论精确解为全 $1$ 向量，数值计算结果也会出现严重偏离。

    \textbf{(3) 考察微小扰动对解的影响}

    扰动实验程序 \texttt{chap-2/Q33.m}：

\begin{lstlisting}[language=Matlab]
% Define the value of n
n = 10;
matrixType = 'polynomial'; % or 'hilbert'

% Generate the original linear system
[A, b] = generateLinearSystem(n, matrixType);

% Exact solution
x_exact = @(n) ones(n, 1);

% Solve the original system
x_original = A \ b;

% Introduce perturbation
A_perturbed = A;
A_perturbed(2, 2) = A_perturbed(2, 2) + 1e-8;
A_perturbed(n, n) = A_perturbed(n, n) + 1e-8;

% Solve the perturbed system
x_perturbed = A_perturbed \ b;

% Compute relative errors
relative_error_exact = norm(x_original - x_exact(n), inf) / norm(x_exact(n), inf);
relative_error_perturbed = norm(x_perturbed - x_exact(n), inf) / norm(x_exact(n), inf);

% Display the results
fprintf('%s\n', matrixType);
fprintf('Relative error of the original solution: %.2e\n', relative_error_exact);
fprintf('Relative error of the perturbed solution: %.2e\n', relative_error_perturbed);
\end{lstlisting}

    分别取 \texttt{matrixType = 'polynomial'} 和 \texttt{matrixType = 'hilbert'} 运行，结果如下：

\begin{lstlisting}
>> Q33
polynomial
Relative error of the original solution: 7.21e-05
Relative error of the perturbed solution: 3.16e-01
>> Q33
hilbert
Relative error of the original solution: 3.18e-04
Relative error of the perturbed solution: 9.77e+00
\end{lstlisting}

    当 $n = 10$ 时，\texttt{polynomial} 和 \texttt{hilbert} 两类线性系统都已表现出明显病态性。此时即使只在矩阵中引入 $10^{-8}$ 量级的微小扰动，也会使解的相对误差显著放大，说明病态方程组对数据误差极其敏感。
\end{solution}

\end{document}
